% BHU PYQ -- Economics: Introductory Macroeconomics (2024-25 NEP)
\documentclass[11pt,a4paper]{article}
\usepackage[a4paper,top=2.5cm,bottom=2.5cm,left=2.8cm,right=2.8cm,headheight=14pt]{geometry}
\usepackage{amsmath,amssymb}
\usepackage{fontspec}
\setmainfont{Kohinoor Devanagari}
\usepackage{microtype,fancyhdr,enumitem,hyperref,lastpage,parskip}

\hypersetup{colorlinks=true,urlcolor=black,linkcolor=black,pdftitle={ECOMJ21: Introductory Macroeconomics -- BHU NEP 2024-25}}

\pagestyle{fancy}\fancyhf{}
\renewcommand{\headrulewidth}{0.4pt}\renewcommand{\footrulewidth}{0.4pt}
\fancyhead[L]{\small \textit{Banaras Hindu University}}
\fancyhead[R]{\small \textit{ECOMJ21: Macroeconomics (2024-25)}}
\fancyfoot[C]{\small Page \thepage\ of \pageref{LastPage}}

\newlist{parts}{enumerate}{2}
\setlist[parts,1]{label=(\alph*),leftmargin=*,itemsep=4pt,topsep=3pt}
\setlist[parts,2]{label=(\roman*),leftmargin=*,itemsep=2pt,topsep=2pt}

\newcommand{\pts}[1]{\hfill \textit{[#1]}}

\begin{document}

\begin{center}
  {\large\bfseries BANARAS HINDU UNIVERSITY}\\[2pt]
  {\normalsize B. A. (Hons.) Semester II Examination, 2024-25 (NEP)}\\[6pt]
  \rule{0.8\linewidth}{0.5pt}\\[6pt]
  {\large\bfseries ECONOMICS}\\[4pt]
  {\large\bfseries Paper Code: ECOMJ-21 : Introductory Macroeconomics (Major)}\\[4pt]
  {\normalsize Time: 2:30 Hours \hfill Full Marks: 70}\\[2pt]
  {\small REF NO. D-II/09}
\end{center}
\rule{\linewidth}{0.4pt}
\smallskip

\noindent \textit{Instructions: This question paper comprises three Sections A, B and C. All questions are compulsory. Write your Roll No.\ at the top immediately on receipt of this question paper.}

\smallskip
\noindent \textit{इस प्रश्न-पत्र में तीन खण्ड अ, ब तथा स हैं। सभी प्रश्न अनिवार्य हैं।}

\bigskip

\section*{SECTION -- A / खण्ड -- अ}
\noindent \textit{Note: Answer each question within 50 words. Each question carries 2 marks.}\pts{5 $\times$ 2 = 10}

\noindent \textit{प्रत्येक प्रश्न का उत्तर 50 शब्दों के भीतर दीजिए। प्रत्येक प्रश्न 2 अंकों का है।}

\begin{enumerate}[label=\textbf{\arabic*.},leftmargin=*,itemsep=8pt]
  \item Explain the following concepts:\pts{5 $\times$ 2 = 10}
  \begin{parts}
    \item Gross Domestic Product (GDP) / सकल घरेलू उत्पाद
    \item Marginal Propensity to Consume (MPC) / सीमांत उपभोग प्रवृत्ति
    \item Inflation Rate / मुद्रास्फीति दर
    \item Money Multiplier / मुद्रा गुणक
    \item High-Powered Money / उच्च-शक्तिशाली मुद्रा
  \end{parts}
\end{enumerate}

\bigskip

\section*{SECTION -- B / खण्ड -- ब}
\noindent \textit{Note: Write your answer in about 250 words. Each question carries 10 marks.}\pts{3 $\times$ 10 = 30}

\noindent \textit{अपना उत्तर लगभग 250 शब्दों में लिखिए। प्रत्येक प्रश्न 10 अंकों का है।}

\begin{enumerate}[label=\textbf{\arabic*.},leftmargin=*,start=2,itemsep=12pt]

  \item Explain the circular flow of income in a three-sector economy (including Households, Firms, and Government).\pts{10}
  
  त्रि-क्षेत्रीय अर्थव्यवस्था (परिवार, फर्म तथा सरकार) में आय के चक्रीय प्रवाह की व्याख्या कीजिए।

  \begin{center}\textbf{OR / अथवा}\end{center}

  Discuss the methods of measuring National Income. What are the major difficulties in measuring National Income in a developing country like India?\pts{10}
  
  राष्ट्रीय आय के मापन की विधियों की चर्चा कीजिए। भारत जैसे विकासशील देश में राष्ट्रीय आय के मापन में आने वाली मुख्य कठिनाइयाँ क्या हैं?

  \item Discuss the Keynesian theory of Income and Employment determination. How is effective demand determined?\pts{10}
  
  आय एवं रोज़गार निर्धारण के केन्ज़ियन सिद्धान्त की विवेचना कीजिए। प्रभावी माँग का निर्धारण किस प्रकार होता है?

  \begin{center}\textbf{OR / अथवा}\end{center}

  Explain the Quantity Theory of Money (Fisher's equation and Cambridge cash balance approach).\pts{10}
  
  मुद्रा के परिमाण सिद्धान्त (फिशर का समीकरण तथा कैम्ब्रिज नकद शेष उपागम) की व्याख्या कीजिए।

  \item Discuss the methods of measuring cost of living in terms of Consumer Price Index (CPI) and GDP Deflator. Highlight their differences, uses and limitations.\pts{10}
  
  उपभोक्ता मूल्य सूचकांक (CPI) तथा GDP अवस्फीतिकारक (Deflator) के पदों में जीवन-निवाह लागत मापने की विधियों की विवेचना कीजिए। उनके अन्तर, उपयोग तथा सीमाओं को रेखांकित कीजिए।

  \begin{center}\textbf{OR / अथवा}\end{center}

  Discuss the interaction between the Multiplier and the Accelerator. How does this interaction explain cyclical fluctuations in income and output?\pts{10}
  
  गुणक तथा त्वरक के मध्य अन्तरक्रिया की विवेचना कीजिए। यह अन्तरक्रिया आय तथा उत्पादन में चक्रीय उच्चावचनों की व्याख्या किस प्रकार करती है?

\end{enumerate}

\bigskip

\section*{SECTION -- C / खण्ड -- स}
\noindent \textit{Note: Write your answer in about 500 words. Each question carries 15 marks.}\pts{2 $\times$ 15 = 30}

\noindent \textit{अपना उत्तर लगभग 500 शब्दों में लिखिए। प्रत्येक प्रश्न 15 अंकों का है।}

\begin{enumerate}[label=\textbf{\arabic*.},leftmargin=*,start=5,itemsep=14pt]

  \item Explain the IS-LM model of macroeconomic equilibrium. How do monetary and fiscal policies shift the IS and LM curves to affect national income and interest rate?\pts{15}
  
  समष्टि अर्थशास्त्र संतुलन के IS-LM मॉडल की व्याख्या कीजिए। मौद्रिक एवं राजकोषीय नीतियाँ राष्ट्रीय आय तथा ब्याज दर को प्रभावित करने के लिए IS तथा LM वक्रों को किस प्रकार विस्थापित करती हैं?

  \begin{center}\textbf{OR / अथवा}\end{center}

  What is Inflation? Discuss its causes, types, and social-economic consequences. Suggest appropriate monetary and fiscal measures to control inflation.\pts{15}
  
  मुद्रास्फीति क्या है? इसके कारणों, प्रकारों तथा सामाजिक-आर्थिक परिणामों की विवेचना कीजिए। मुद्रास्फीति को नियंत्रित करने के लिए उपयुक्त मौद्रिक एवं राजकोषीय उपायों का सुझाव दीजिए।

  \item Discuss the Classical theory of Employment and Output. How did Keynes criticize Say's Law of Markets and Pigou's wage-cut policy?\pts{15}
  
  रोज़गार एवं उत्पादन के प्रतिष्ठित (क्लासिकल) सिद्धान्त की विवेचना कीजिए। केन्ज़ ने से (Say) के बाजार नियम तथा पीगू की मजदूरी-कटौती नीति की आलोचना किस प्रकार की?

  \begin{center}\textbf{OR / अथवा}\end{center}

  Critically examine the Consumption Function hypothesis propounded by Keynes. Explain the Absolute Income Hypothesis and Relative Income Hypothesis.\pts{15}
  
  केन्ज़ द्वारा प्रतिपादित उपभोग फलन उपकल्पना का समालोचनात्मक परीक्षण कीजिए। निरपेक्ष आय उपकल्पना तथा सापेक्ष आय उपकल्पना की व्याख्या कीजिए।

\end{enumerate}

\end{document}
